% ch10_supplement.tex — 第10章：等离子体约束与受控核聚变 补充内容
% ============================================================
% 经典例题 + 完整解答、TikZ图像、核心公式速查卡、自学检查清单

% ============================================================
% 1. 经典例题
% ============================================================

\section{经典例题与解答}

\begin{example}{Lawson判据的数值验证}{ex:lawson}
对于氘-氚（D-T）聚变反应，在温度 $T=10\,\mathrm{keV}$ 时，验证Lawson判据并求所需的最小 $n\tau$ 值。
\end{example}

\textbf{已知条件：}
\begin{itemize}
    \item D-T反应截面参数：在 $T=10\,\mathrm{keV}$ 时，$\langle\sigma v\rangle_{DT} \approx 1.1\times10^{-22}\,\mathrm{m^3\cdot s^{-1}}$
    \item D-T反应每次释放能量：$E_{\alpha}=3.5\,\mathrm{MeV}$（$\alpha$粒子），$E_n=14.1\,\mathrm{MeV}$（中子）
    \item 考虑能量约束效率 $\eta=0.3$（仅$\alpha$粒子能量加热等离子体）
\end{itemize}

\textbf{解答：}

Lawson判据的核心思想是：聚变反应产生的\textbf{可自持}功率必须大于等于等离子体损失的能量功率。先写\emph{能量收支}，再谈门槛。

\textbf{能量收支表（全体积功率密度，D-T 等量、$n\equiv n_e=n_D+n_T$、$T_e=T_i=T$）：}
\begin{itemize}
    \item 等离子体内能密度：$u=3nk_BT$（电子与离子各 $\tfrac32nk_BT$）
    \item $\alpha$ 粒子自加热：$P_\alpha=\dfrac{n^2}{4}\langle\sigma v\rangle E_\alpha$（$\tfrac{n}{2}$ 个 D 与 $\tfrac{n}{2}$ 个 T，反应率 $\propto \tfrac{n^2}{4}\langle\sigma v\rangle$；仅 $\alpha$ 留在等离子体中）
    \item 韧致辐射损失：$P_{\rm rad}=C_B n^2\sqrt{T_e}$（$T_e$ 以 eV 计，$C_B\approx1.69\times10^{-38}\,\mathrm{W\,m^3\,eV^{-1/2}}$，$Z_{\rm eff}=1$）
    \item 由输运造成的能量损失：$P_{\rm loss}=\dfrac{3nk_BT}{\tau_E}$（$\tau_E$ 定义为\textbf{非辐射}输运约束时间）
\end{itemize}
点火条件是 $P_\alpha\ge P_{\rm rad}+P_{\rm loss}$，即
\begin{equation}
    \frac{n^2}{4}\langle\sigma v\rangle E_\alpha \ge C_B n^2\sqrt{T_e} + \frac{3nk_BT}{\tau_E}
\label{eq:ch10-lawson-full}
\end{equation}
解出 $n\tau_E$（注意 $\tfrac{P_{\rm rad}}{n^2}$ 与 $\tfrac{P_\alpha}{n^2}$ 都是\textbf{单位 $n^2$ 的功率密度}，量纲一致；\textbf{不要再乘一个 $n$}）：
\begin{equation}
    \boxed{\;n\tau_E \ge \frac{3k_BT}{\dfrac{E_\alpha}{4}\langle\sigma v\rangle - \dfrac{P_{\rm rad}}{n^2}}\;}
\label{eq:ch10-lawson-full-solved}
\end{equation}
\textbf{分母必须为正}——这本身就是一条物理要求：若 $\tfrac{E_\alpha}{4}\langle\sigma v\rangle<\tfrac{P_{\rm rad}}{n^2}$，$\alpha$ 加热根本压不住辐射，无论 $n\tau_E$ 多大都无法点火（这就是存在"最低点火温度"的原因，约 $4\,\mathrm{keV}$ 以下 D-T 无法点火）。

\textbf{忽略辐射时的简化式：} 令 $P_{\rm rad}\to0$，得
\begin{equation}
    n\tau_E \ge \frac{12k_B T}{E_{\alpha}\langle\sigma v\rangle}
\label{eq:ch10-lawson-simplified}
\end{equation}

\textbf{(1) 忽略辐射：} 注意单位换算 $1\,\mathrm{keV}=1.602\times10^{-16}\,\mathrm{J}$：
\[
    12k_BT = 12\times1.602\times10^{-16}\times10 = 1.9224\times10^{-14}\,\mathrm{J}
\]
\[
    E_\alpha\langle\sigma v\rangle = 3.5\times10^3\times1.602\times10^{-16}\times1.1\times10^{-22}
    = 6.1677\times10^{-35}\,\mathrm{J\,m^3\,s^{-1}}
\]
\[
    n\tau_E \ge \frac{1.9224\times10^{-14}}{6.1677\times10^{-35}} = 3.12\times10^{20}\,\mathrm{m^{-3}\,s}
\]

\textbf{(2) 加入辐射：门槛\emph{升高}，不是降低。} 由 (1) 的系数算辐射项（与 $n$ 无关）：
\[
    \frac{P_{\rm rad}}{n^2} = C_B\sqrt{T_e} = 1.69\times10^{-38}\times\sqrt{10^4} = 1.69\times10^{-36}\,\mathrm{J\,m^3\,s^{-1}}
\]
\[
    \frac{E_\alpha}{4}\langle\sigma v\rangle = \frac{6.1677\times10^{-35}}{4} = 1.542\times10^{-35}
\]
分母：
\[
    1.542\times10^{-35} - 1.69\times10^{-36} = 1.373\times10^{-35}
\]
\[
    n\tau_E \ge \frac{3\times1.602\times10^{-16}\times10}{1.373\times10^{-35}}
    = \frac{4.806\times10^{-15}}{1.373\times10^{-35}} \approx \boxed{3.5\times10^{20}\,\mathrm{m^{-3}\,s}}
\]

\textbf{核对方向：} 忽略辐射给 $3.12\times10^{20}$，加入辐射后升到 $3.5\times10^{20}$。加入一项\emph{损失}\textbf{只能让门槛更高}——把门槛从 $3.12$ 降到 $2$（原稿的错误结论）在能量收支上是反向的。

\textbf{与文献值的差别来自模型选择，不是矛盾：} 常见的"$n\tau_E\approx1.5\text{--}2\times10^{20}$"是\emph{不同约定}下的数（例如把 $n$ 取为纯 D-T 燃料密度而非 $n_e$、或只要求 $P_\alpha\ge P_{\rm loss}$ 而把辐射单列、或 $\eta$ 的吸收因子分摊方式不同），模型一变阈值就变。\textbf{必须连同模型假设一起引用}，不能把不同约定下的数当成同一个"标准值"混比。此外，$\eta=0.3$ 在本模型里对应"$\alpha$ 能量的 30\% 真正沉积"这一修正——若采用它，应把 $P_\alpha$ 乘以 0.3（阈值升高约 3.3 倍），\textbf{不能列出后不使用}。

\textbf{温度依赖：} 随温度升高，$\langle\sigma v\rangle$ 增大而 $\sqrt{T_e}$ 增大得慢，故分母中的"净加热"变强，$n\tau_E$ 门槛下降；D-T 在 $15\text{--}25\,\mathrm{keV}$ 附近门槛最低（约 $1.5\times10^{20}\,\mathrm{m^{-3}s}$ 量级），温度更高时因韧致辐射增加又回升。

% ------------------------------------------------------------------------

\begin{example}{三重乘积 $nT\tau_E$ 计算}{ex:triple}
某托卡马克装置参数如下：
\begin{itemize}
    \item 大半径 $R=1.7\,\mathrm{m}$，小半径 $a=0.56\,\mathrm{m}$
    \item 中心磁场 $B_0=3.0\,\mathrm{T}$
    \item 平均电子密度 $n_e=5\times10^{19}\,\mathrm{m^{-3}}$
    \item 电子温度 $T_e=8\,\mathrm{keV}$（假设 $T_i=T_e$）
    \item 能量约束时间 $\tau_E=0.5\,\mathrm{s}$
\end{itemize}
计算该装置的三重乘积 $nT\tau_E$，并与D-T点火所需的典型值 $nT\tau_E\approx 3\times10^{21}\,\mathrm{keV\cdot m^{-3}\cdot s}$ 比较。
\end{example}

\textbf{解答：}

三重乘积（Triple Product）是衡量磁约束聚变装置性能的关键指标，定义为：
\begin{equation}
    nT\tau_E = n_e \times T_e \times \tau_E
\label{eq:ch10-triple-product}
\end{equation}

（注：离子温度 $T_i$ 与电子温度 $T_e$ 近似相等时，可直接使用 $T=T_e$。）

将给定数值代入：
\[
    nT\tau_E = (5\times10^{19}\,\mathrm{m^{-3}}) \times (8\,\mathrm{keV}) \times (0.5\,\mathrm{s})
\]
\[
    = 5\times8\times0.5 \times 10^{19}\,\mathrm{keV\cdot m^{-3}\cdot s}
\]
\[
    = 20 \times 10^{19}\,\mathrm{keV\cdot m^{-3}\cdot s}
\]
\[
    \boxed{nT\tau_E = 2\times10^{20}\,\mathrm{keV\cdot m^{-3}\cdot s}}
\]

\textbf{与点火条件比较：}

D-T反应的点火条件（考虑了能量自持和辐射损失）通常要求：
\[
    nT\tau_E \gtrsim 3\times10^{21}\,\mathrm{keV\cdot m^{-3}\cdot s}
\]

该装置的三重乘积为 $2\times10^{20}$，约为点火所需值的 $7\%$。因此，该装置尚远未达到点火条件。

\textbf{改进方向分析：}
\begin{itemize}
    \item 若将密度提升至 $n_e=1\times10^{20}\,\mathrm{m^{-3}}$，温度提升至 $T_e=15\,\mathrm{keV}$，约束时间 $\tau_E=2\,\mathrm{s}$：
    \[
        nT\tau_E = 1\times10^{20}\times15\times2 = 3\times10^{21}\,\mathrm{keV\cdot m^{-3}\cdot s}
    \]
    此时接近点火条件。
\end{itemize}

% ------------------------------------------------------------------------

\begin{example}{托卡马克与仿星器的约束时间对比}{ex:conftime}
比较以下两种装置的约束时间特性：
\begin{itemize}
    \item 托卡马克A：大半径 $R=2.0\,\mathrm{m}$，小半径 $a=0.6\,\mathrm{m}$，环向磁场 $B=4.0\,\mathrm{T}$，等离子体电流 $I_p=5\,\mathrm{MA}$
    \item 仿星器B：大半径 $R=2.0\,\mathrm{m}$，平均小半径 $a=0.5\,\mathrm{m}$，环向磁场 $B=2.5\,\mathrm{T}$，旋转变换 $\iota/2\pi = 0.5$
\end{itemize}
利用经验定标律估算各自的能量约束时间。
\end{example}

\textbf{解答：}

\textbf{模型说明（先读）：} 下面的定标式都是\emph{教学用简化式}，不是真实装置的定标律。为避免"算出一个数后又挑一个更顺眼的数替代"，这里\textbf{只用一个假设的定标式比较两类装置}，并把计算结果与模型偏差一起保留。真实托卡马克与仿星器有各自的定标数据库（ITER89-P/98y2、ISS04 等），幂次、单位、运行模式与拟合范围都不同，不能由这一例推出"托卡马克一定优于仿星器"。

\textbf{统一假设的定标式（本书教学用，非真实定标）：}
\begin{equation}
    \tau_E^{\text{(教)}} = 0.04 \cdot H \cdot I_p^{1.25} \cdot B^{0.5} \cdot n_{19}^{0.25} \cdot a^{1.5} \cdot R^{1.75}\quad\text{(托卡马克，Goldston L-模量级)}
    \label{eq:ch10-tokamak-scaling}
\end{equation}
（单位：$I_p$ 为 MA，$B$ 为 T，$n_{19}$ 为 $10^{19}\,\mathrm{m^{-3}}$，$a,R$ 为 m，$\tau_E$ 为 s。此式的幂次取自 Goldston L-模，用于教学量级估算。）

假设托卡马克A的 $n=5\times10^{19}\,\mathrm{m^{-3}}$（即 $n_{19}=5$），$H=1$（L-模）：
\[
    \tau_E^{Tok} = 0.04 \times 1 \times 5^{1.25} \times 4^{0.5} \times 5^{0.25} \times 0.6^{1.5} \times 2.0^{1.75}
\]

逐步计算：
\begin{align*}
    5^{1.25} &= 5 \times 5^{0.25} \approx 5 \times 1.495 = 7.475 \\
    4^{0.5} &= 2 \\
    5^{0.25} &\approx 1.495 \\
    0.6^{1.5} &= 0.6\times\sqrt{0.6} \approx 0.6\times0.775 = 0.465 \\
    2.0^{1.75} &= 2^1 \times 2^{0.75} = 2 \times 1.682 = 3.364
\end{align*}

\[
    \tau_E^{Tok} = 0.04 \times 7.475 \times 2 \times 1.495 \times 0.465 \times 3.364
\]
\[
    = 0.04 \times 34.95 \approx 1.40\,\mathrm{s}
\]

\textbf{仿星器部分：不用"另一条定标律"，只做幂律比较。} 仿星器没有等离子体电流，故 (上面的) $I_p^{1.25}$ 项不适用；它与托卡马克最本质的差别在于：约束不再由电流驱动，而由外线圈的旋转变换决定。若继续把同一条教学幂律套到仿星器B上（把 $I_p$ 当作一个等效量、$a=0.5$、$B=2.5$、$R=2.0$、$n_{19}=3$），会得到
\[
    \tau_E^{Stel} \approx 0.078\,\mathrm{s}
\]
\textbf{但这个数不能当真}——式中的 $I_p$ 幂次根本没有物理依据。\textbf{正确的做法是不要把它当结果报出}，而是明确：\underline{本教学模型不含仿星器的正确标度，故不给仿星器的定量结论}。真实仿星器（如 W7-X，$R\approx5.5\,\mathrm{m}$）在优化位形下同样能达到秒量级约束，与托卡马克并无"必然更差"的排序。

\textbf{结论：}
\begin{itemize}
    \item 托卡马克A：在本教学定标式下 $\tau_E \approx 1.4\,\mathrm{s}$（电流驱动的极向场使约束时间随 $I_p$ 显著提高）。
    \item 仿星器B：\textbf{本模型不适用}，不给定量值。两类装置的定量比较需要各自的专用定标律与运行模式。
    \item 物理差别是可靠的：托卡马克靠等离子体电流获得极向场，存在破裂风险；仿星器用外线圈产生旋转变换，可稳态运行。这是两种设计路线的\emph{定性}区别，与上面的数字无关。
\end{itemize}

% ------------------------------------------------------------------------

\begin{example}{ITER参数估算}{ex:iter}
ITER（国际热核聚变实验堆）主要设计参数如下：
\begin{itemize}
    \item 大半径 $R=6.2\,\mathrm{m}$
    \item 小半径 $a=2.0\,\mathrm{m}$
    \item 环向磁场（轴上）$B_0=5.3\,\mathrm{T}$
    \item 等离子体电流 $I_p=15\,\mathrm{MA}$
    \item 拉长比 $\kappa = 1.7$
    \item 三角形变 $\delta = 0.33$
\end{itemize}
请估算：
\begin{enumerate}
    \item 等离子体体积 $V_p$
    \item 安全因子（q-因子）在边缘处的近似值 $q_{95}$
\end{enumerate}
\end{example}

\textbf{解答：}

\textbf{（1）等离子体体积 $V_p$}

对于非圆截面托卡马克，等离子体体积可近似为环形体积乘以截面修正：
\begin{equation}
    V_p \approx 2\pi R \times \pi a^2 \times \kappa
\label{eq:ch10-tokamak-volume}
\end{equation}

其中 $2\pi R$ 为环向大圆周长，$\pi a^2$ 为圆形截面面积，$\kappa$ 为垂直方向拉长比。

代入数值：
\[
    V_p = 2\pi \times 6.2 \times \pi \times (2.0)^2 \times 1.7
\]
\[
    = 2\pi \times 6.2 \times 4\pi \times 1.7
\]
\[
    = 2 \times 6.2 \times 4 \times 1.7 \times \pi^2
\]
\[
    = 84.32 \times 9.87
\]
\[
    \boxed{V_p \approx 832\,\mathrm{m^3}}
\]

（ITER官方设计值为 $V_p \approx 840\,\mathrm{m^3}$，此估算相当精确。）

\textbf{（2）安全因子 $q_{95}$}

安全因子 $q$ 是托卡马克中重要的稳定性参数，定义为磁场线在极向转一圈时环向转过的圈数：
\begin{equation}
    q = \frac{r B_{\phi}}{R B_{\theta}}
\label{eq:ch10-safety-factor}
\end{equation}

在边缘处（$r=a$），利用安培定律估算极向磁场 $B_{\theta}$：
\begin{equation}
    \oint \vect{B}_{\theta}\cdot\diff\vect{l} = \mu_0 I_{enc} \quad\Rightarrow\quad 2\pi a B_{\theta} = \mu_0 I_p
\label{eq:ch10-ampere-law}
\end{equation}
\begin{equation}
    B_{\theta}(a) = \frac{\mu_0 I_p}{2\pi a}
\label{eq:ch10-poloidal-field}
\end{equation}

代入数值（$I_p=15\,\mathrm{MA}=15\times10^6\,\mathrm{A}$）：
\[
    B_{\theta}(a) = \frac{4\pi\times10^{-7}\times15\times10^6}{2\pi\times2.0}
\]
\[
    = \frac{2\times10^{-7}\times15\times10^6}{2.0}
\]
\[
    = \frac{3.0}{2.0} = 1.5\,\mathrm{T}
\]

边缘安全因子（大柱面近似）：
\[
    q_a = \frac{a B_{\phi}}{R B_{\theta}} = \frac{2.0\times5.3}{6.2\times1.5} = \frac{10.6}{9.3} \approx 1.14
\]

考虑非圆截面修正（通常采用 $q_{95}$，即磁面 $95\%$ 通量处的安全因子）：
\begin{equation}
    q_{95} \approx q_a \times \frac{(1+\kappa^2)}{2} \times \text{三角修正}
\label{eq:ch10-q95-approx}
\end{equation}

对于标准ITER参数，经验公式为：
\begin{equation}
    q_{95} \approx \frac{5 a^2 B_0}{R I_p} \cdot \frac{(1+\kappa^2)}{2}
\label{eq:ch10-q95-empirical}
\end{equation}

代入：
\[
    q_{95} \approx \frac{5\times(2.0)^2\times5.3}{6.2\times15} \times \frac{1+(1.7)^2}{2}
\]
\[
    = \frac{5\times4\times5.3}{93} \times \frac{3.89}{2}
\]
\[
    = \frac{106}{93} \times 1.945
\]
\[
    \approx 1.14 \times 1.945 \approx 2.2
\]

这一简化公式给出的 $q_{95}\approx 2.2$ 量级正确。实际的 ITER 设计目标为 $q_{95}\approx 3.0$，差异来源于更完整的定标公式还包含三角形变、内感等几何修正因子。采用 ITER 官方设计值：
\[
    \boxed{q_{95} \approx 3.0\ \text{（ITER 官方设计值，简化公式估算为 2.2）}}
\]

% ============================================================
% 2. TikZ图像
% ============================================================

% ch10 新例题：惯性约束聚变（ICF）的能量估算
\begin{example}{惯性约束聚变的压缩与点火估算}{icf-estimate}
惯性约束聚变（ICF）靶丸：初始半径 $r_0 = 1\,\mathrm{mm}$，DT 冰层质量 $m = 0.3\,\mathrm{mg}$。
（1）若压缩到初始体积的 $1/1000$，求压缩后的半径与密度（初始 DT 密度 $\rho_0 = 0.25\,\mathrm{g/cm^3}$）；
（2）估算加热到 $T=10\,\mathrm{keV}$ 需要的燃料内能与驱动激光能量（假设流体力学效率 $\eta_h = 10\%$）；
（3）说明惯性约束与磁约束在约束时间上的本质区别。
\tcblower
\textbf{解答：}

\textbf{(1) 压缩后的半径与密度：} 体积压缩 1000 倍对应半径压缩 $\sqrt[3]{1000}=10$ 倍：
\[
r = \frac{r_0}{10} = 0.1\,\mathrm{mm} = 100\,\mu\mathrm{m}, \qquad \rho = 1000\,\rho_0 = 250\,\mathrm{g/cm^3} = 2.5\times10^5\,\mathrm{kg/m^3}
\]
密度 $250\,\mathrm{g/cm^3}$ 约为固体密度的 1000 倍——这正是 ICF"把燃料压缩到极端密度"的含义。

\textbf{几何假设（先声明）：} 为把质量与几何分开，下面\textbf{只用给定的质量 $m=0.3\,\mathrm{mg}$} 与压缩后密度算体积，不把外半径当实心球（外半径与冰层体积本就不是一个均匀实心球）。
\textbf{(2) 燃料内能与驱动激光：} 压缩后的粒子数密度
\[
n = \frac{\rho}{m_i} = \frac{2.5\times10^{5}}{4.15\times10^{-27}} \approx 6.0\times10^{31}\,\mathrm{m^{-3}}
\]
（与直接 $n=\rho/m_i$ 一致。）加热到 $T=10\,\mathrm{keV}$，电子与离子的总内能密度（每自由度 $k_{\!B}T/2$，3 个自由度 $\times$ 2 类粒子，故系数为 $3nk_{\!B}T$）：
\[
u = 3\,n\,k_{\!B}T = 3\times6.0\times10^{31}\times1.602\times10^{-15}
\approx 2.9\times10^{17}\,\mathrm{J/m^3}
\]
\textbf{注意指数：是 $2.9\times10^{17}$，不是 $2.9\times10^{18}$。} 压缩后的燃料体积由给定的质量与密度算得：
\[
V = \frac{m}{\rho} = \frac{0.3\times10^{-6}}{2.5\times10^{5}} = 1.2\times10^{-12}\,\mathrm{m^3}
\]
故燃料内能：
\[
E_{\text{fuel}} = uV \approx 2.9\times10^{17}\times1.2\times10^{-12} \approx 3.5\times10^{-1}\,\mathrm{MJ} = 0.35\,\mathrm{MJ}
\]
\textbf{独立核对（推荐做法）：} 内能也可以用质量直接写——$E_{\text{fuel}}=3(m/m_i)k_{\!B}T=3\times(0.3\times10^{-6}/4.15\times10^{-27})\times1.602\times10^{-15}\approx0.35\,\mathrm{MJ}$，与上式吻合。\textbf{凡是密度与体积都要用的算例，务必用这样的第二条路径复核。}

考虑流体力学效率（激光能量只有 $10\%$ 传给燃料）：
\[
E_{\text{laser}} = \frac{0.35\,\mathrm{MJ}}{0.1} \approx 3.5\,\mathrm{MJ}
\]
这比实际 ICF 装置（NIF 单发到靶 $\approx1.8$--$2.05\,\mathrm{MJ}$）的驱动能量同一量级、略偏大——简化估算用的是"整块燃料均匀加热到 $10\,\mathrm{keV}$"，而真实的中心热斑模型只需点燃一小部分燃料、其余靠 $\alpha$ 自加热与激波压缩，故把需求高估数倍是预期的。\textbf{量级判断合理}；但题目算出的 $35\,\mathrm{MJ}$（原稿）是把 $u$ 的指数多写了一位所致，必须改正。

\textbf{(3) 约束时间的本质区别：} 惯性约束没有外场约束，靠燃料自身的\textbf{惯性}维持：约束时间 $\tau \sim r/c_s \approx 100\,\mu\mathrm{m}/(10^6\,\mathrm{m/s}) \sim 100\,\mathrm{ps}$——极短。磁约束则用磁场把低密度（$10^{20}\,\mathrm{m^{-3}}$）等离子体约束数秒；惯性约束用极高密度（$\sim10^{31}\,\mathrm{m^{-3}}$，固体千倍）把反应时间压缩到纳秒以下。Lawson 判据 $n\tau$ 的两种满足方式：磁约束 $n$ 低 $\tau$ 长，惯性约束 $n$ 极高 $\tau$ 极短——这正是两种约束策略的物理根源。
\end{example}

\section{关键装置与参数空间图}

% ------------------------------------------------------------------------
% 图1：托卡马克位形示意图
% ------------------------------------------------------------------------
\begin{figure}[htbp]
    \centering
    \begin{tikzpicture}[scale=1.0,>=Stealth]
        % 坐标轴
        \draw[->] (-0.5,0) -- (7,0) node[right] {$R$};
        \draw[->] (0,-3.5) -- (0,3.5) node[above] {$Z$};
        
        % 环向轴（中心轴）
        \draw[dashed, gray] (2,0) -- (2,3);
        \draw[dashed, gray] (2,0) -- (2,-3);
        \node[gray] at (2,-3.3) {环向轴};
        
        % 等离子体截面（拉长椭圆）
        \draw[thick, blue, fill=blue!10] (4.5,0) ellipse (1.5 and 2.5);
        \node[blue] at (4.2,-1.6) {\small 等离子体};

        % 大半径 R
        \draw[<->, thick, red] (2,2.8) -- (4.5,2.8);
        \node[red, above] at (3.25,2.9) {$R$};

        % 小半径 a
        \draw[<->, thick, red] (5.2,-2.5) -- (5.2,2.5);
        \node[red, right] at (5.3,0) {$a$};

        % 环向磁场 B_phi（虚线大圆，俯视图示意）
        \draw[dashed, gray] (2,0) circle (2.5);
        \draw[->, thick, green!60!black] (4.5,0) arc (0:30:2.5);
        \node[green!60!black] at (5.6,1.3) {$B_{\phi}$};

        % 极向磁场 B_theta（截面内圆圈箭头）
        \draw[->, thick, orange] (3.8,0) arc (180:90:0.7);
        \node[orange] at (3.0,0.8) {$B_{\theta}$};

        % 等离子体电流 I_p
        \draw[->, ultra thick, purple] (4.9,-0.5) -- (4.9,0.1);
        \node[purple, right] at (4.95,0.35) {$I_p$};
        
        % 图例说明
        \node[anchor=west] at (6.5,2) {\color{green!60!black}\small 环向磁场：};
        \node[anchor=west] at (6.5,1.5) {\color{green!60!black}\small 外部线圈产生};
        
        \node[anchor=west] at (6.5,0.5) {\color{orange}\small 极向磁场：};
        \node[anchor=west] at (6.5,0) {\color{orange}\small 等离子体电流产生};
        
        \node[anchor=west] at (6.5,-1) {\color{purple}\small $I_p$：等离子体电流};
    \end{tikzpicture}
    \caption{托卡马克位形示意图（$R$为大半径，$a$为小半径，$B_{\phi}$为环向磁场，$B_{\theta}$为极向磁场，$I_p$为等离子体电流）。}
    \label{fig:tokamak}
\end{figure}

% ------------------------------------------------------------------------
% 图2：Lawson参数空间图
% ------------------------------------------------------------------------
\begin{figure}[htbp]
    \centering
    \begin{tikzpicture}[scale=1.0,>=Stealth]
        % 坐标轴
        \draw[->] (0,0) -- (10,0) node[right] {$T\,\mathrm{(keV)}$};
        \draw[->] (0,0) -- (0,6) node[above] {$\log_{10}(n\tau)\;\mathrm{(m^{-3}\cdot s)}$};
        
        % 刻度
        \foreach \x in {1,2,3,4,5,6,7,8,9} {
            \draw (\x,0.1) -- (\x,-0.1) node[below] {\x};
        }
        \foreach \y in {19,20,21,22} {
            \pgfmathsetmacro{\yp}{\y-18}
            \draw (0.1,\yp) -- (-0.1,\yp) node[left] {\y};
        }
        
        % D-T 点火线（示意曲线）
        \draw[thick, red, domain=1.5:9, samples=50] plot (\x, {1+4/(\x^0.5)});
        \node[red] at (7.5,2.6) {D-T 点火线};
        \node[red] at (7.5,2.2) {\small $n\tau \sim T^{1/2}$};

        % D-D 点火线（更高要求）
        \draw[thick, blue, dashed, domain=2:9, samples=50] plot (\x, {2+4/(\x^0.5)});
        \node[blue] at (8.4,3.3) {D-D 点火线};

        % 装置标注
        % TFTR
        \filldraw[black] (3.5,2.5) circle (3pt);
        \node[below] at (3.5,2.35) {TFTR};

        % JET
        \filldraw[black] (5,3.2) circle (3pt);
        \node[above right] at (5,3.2) {JET};

        % DIII-D
        \filldraw[black] (4,2.8) circle (3pt);
        \node[below right] at (4,2.8) {DIII-D};

        % ITER（设计目标）
        \filldraw[red] (8,4.35) circle (4pt);
        \node[above right, red] at (8,4.35) {ITER（设计目标）};

        % NIF（惯性约束，不同的n-tau区域，这里大致标出）
        \filldraw[blue] (2.5,3.8) circle (3pt);
        \node[above left, blue] at (2.5,3.8) {NIF};

        % 现有装置区域
        \draw[dotted, gray, thick] (2,2) -- (6,2.5) -- (6.5,3.5) -- (2,3) -- cycle;
        \node[gray] at (2.9,2.6) {\small 现有磁约束};
        \node[gray] at (2.9,2.3) {\small 装置区域};

        % 点火区域标注：在点火线【上方】（相同 T 下 n tau 更高才满足）
        \draw[red, fill=red!10, opacity=0.3] (5,3.5) -- (9,3.5) -- (9,4.5) -- (5,4.5) -- cycle;
        \node[red, opacity=0.7] at (6.3,4.0) {点火区（阈值线\emph{上方}）};

        % 三个判据状态点（固定 T，看 n tau 是否越过门槛）
        \filldraw[red] (3.0,1.9) circle (2.5pt);
        \node[red,right,font=\footnotesize] at (3.05,1.75) {\cnum{1}不满足};
        \filldraw[orange] (3.0,2.55) circle (2.5pt);
        \node[orange,right,font=\footnotesize] at (3.05,2.4) {\cnum{2}临界};
        \filldraw[green!50!black] (3.0,3.2) circle (2.5pt);
        \node[green!50!black,right,font=\footnotesize] at (3.05,3.05) {\cnum{3}满足};
    \end{tikzpicture}
    \caption{Lawson参数空间图：$n\tau$ 对温度 $T$ 的依赖关系。图中标注了若干代表性装置（TFTR、JET、DIII-D、NIF）以及ITER的设计目标。\textbf{判据方向：固定温度时，必须 $n\tau$ 高于点火线才满足点火条件，所以能量自持区在点火线\emph{上方}（图中红色阴影区）；线下方为不满足区。}三点示意：\cnum{1}在阈值线下（不满足）、\cnum{2}恰在线上（临界）、\cnum{3}在线上方（满足）。}
    \label{fig:lawson}
\end{figure}

% ============================================================
% 3. 核心公式速查卡
% ============================================================

\section{核心公式速查卡}

\begin{formulabox}{Lawson判据与三重乘积}{f:lawson}
对于D-T反应，Lawson判据（温度形式）：
\begin{equation}
    n\tau \geq \frac{12k_B T}{E_{\alpha}\langle\sigma v\rangle}
\label{eq:ch10-lawson-simplified-ex}
\end{equation}

工程上常用的三重乘积：
\[
    nT\tau_E \geq 3\times10^{21}\,\mathrm{keV\cdot m^{-3}\cdot s}\quad\text{(D-T点火)}
\]

\textbf{常用反应速率参数}（$T$ 单位为keV）：
\begin{align*}
    \langle\sigma v\rangle_{DT} &\approx 1.1\times10^{-24}\,T^2\quad (T\approx10\sim20\,\mathrm{keV})\\
    \langle\sigma v\rangle_{DD} &\approx 3\times10^{-25}\quad (T\approx50\,\mathrm{keV})
\end{align*}
\end{formulabox}

\begin{formulabox}{托卡马克基本参数}{f:tokamak}
等离子体体积（拉长截面）：
\[
    V_p \approx 2\pi^2 R a^2 \kappa
\]

安全因子（圆柱近似）：
\[
    q(r) = \frac{r B_{\phi}}{R B_{\theta}} \quad\Rightarrow\quad q_a = \frac{a B_0}{R B_{\theta}(a)}
\]

边缘极向磁场（安培定律）：
\begin{equation}
    B_{\theta}(a) = \frac{\mu_0 I_p}{2\pi a}
\label{eq:ch10-poloidal-field-ex}
\end{equation}

$q_{95}$ 经验公式：
\[
    q_{95} \approx \frac{5 a^2 B_0}{R I_p}\cdot\frac{(1+\kappa^2)}{2}
\]
（$a$ 单位为m，$B_0$ 为T，$R$ 为m，$I_p$ 为MA）
\end{formulabox}

\begin{formulabox}{磁约束约束时间经验定标律}{f:scaling}
\textbf{托卡马克（Goldston L-模）}：
\[
    \tau_E = 0.04 \cdot H \cdot I_p^{1.25} \cdot B^{0.5} \cdot n_{19}^{0.25} \cdot a^{1.5} \cdot R^{1.75}
\]

\textbf{托卡马克（ITER98y2 H-模）}：
\[
    \tau_E = 0.0562 \cdot H \cdot I_p^{0.93} \cdot B^{0.15} \cdot n_{19}^{0.41} \cdot a^{1.97} \cdot R^{1.39} \cdot M^{0.19} \cdot P^{-0.58}
\]

\textbf{仿星器（W7-AS类型）}：
\[
    \tau_E \propto a^2 R^{0.5} B^{0.5} n^{0.3}
\]
\end{formulabox}

\begin{formulabox}{韧致辐射功率损失}{f:brems}
电子韧致辐射（非相对论、光学薄、纯氢 $Z=1$ 近似）：
\[
    P_{rad} = \frac{2^{5/2}\pi^{1/2}e^6}{3(4\pi\eps_0)^3 m_e^{3/2}c^3\hbar} \cdot \bar{g} \cdot n_e^2 \sqrt{T_e}
\]
其中 $\bar g$ 是热平均冈特因子（量级 1--3，氢等离子体约 $2.5$）。

简化数值形式（$T_e$ 以 $\mathrm{eV}$ 计，已含 $\bar g\approx2.5$；与本章 Lawson 例题所用的 $C_B$ 一致）：
\[
    P_{rad}\;[\mathrm{W\cdot m^{-3}}] \approx 1.69\times10^{-38} \cdot n_e^2\;[\mathrm{m^{-6}}] \cdot \sqrt{T_e\;[\mathrm{eV}]}
\]
\textbf{注意系数口径：} 若只取裸的非相对论分子（$\bar g=1$），系数约为 $6.6\times10^{-39}$；文献中另见 $5.3\times10^{-37}$ 一类的值，那是含相对论修正与频率积分的韧致辐射连续谱总量，口径与上式不同。\textbf{同一本书内必须用同一个 $C_B$}——本章的能量收支与 Lawson 推演一律采用带 $\bar g\approx2.5$ 的 $1.69\times10^{-38}$。
\end{formulabox}

% ============================================================
% 4. 自学检查清单
% ============================================================

\section{自学检查清单}

\begin{checklistbox}{第10章自学检查清单}{ch10-checklist}
\begin{itemize}[leftmargin=*,label=$\square$]
\item[$\square$] 我能写出D-T和D-D聚变反应方程式，并说明它们各自的能量释放特点
\item[$\square$] 我理解Lawson判据的物理含义：为什么需要密度和约束时间的乘积？
\item[$\square$] 我能从能量平衡推导 $n\tau$ 的表达式，并解释辐射损失如何修正判据
\item[$\square$] 我了解三重乘积 $nT\tau_E$ 与Lawson判据的关系，以及工程上的点火标准
\item[$\square$] 我能画出托卡马克的截面图，标出 $R$、$a$、$B_{\phi}$、$B_{\theta}$、$I_p$
\item[$\square$] 我理解安全因子 $q$ 的物理意义：它如何衡量磁场线的螺旋程度？
\item[$\square$] 我能用安培定律估算托卡马克边缘的极向磁场 $B_{\theta}(a)$
\item[$\square$] 我了解托卡马克与仿星器在磁场位形产生方式上的根本区别
\item[$\square$] 我知道惯性约束聚变（ICF）与磁约束聚变（MCF）在约束原理上的不同
\item[$\square$] 我能列举至少三个现有聚变装置（如JET、W7-X、EAST、NIF）及其主要特点
\item[$\square$] 我理解ITER的设计目标以及为什么需要如此巨大的尺寸
\item[$\square$] 我能解释为什么 $q=1$ 和 $q=2$ 面对托卡马克 MHD 稳定性有特殊意义
\end{itemize}
\end{checklistbox}

\endinput
